%%%%%%%%%%%%%%%%%%%%%%%%%%%%%%%%%%%%%%%%%%%%%%%%%%%%%%%%%%%%%%%%%%%%%%%%
\section{Transfer to finite meanders}\label{sec:bridge}
%%%%%%%%%%%%%%%%%%%%%%%%%%%%%%%%%%%%%%%%%%%%%%%%%%%%%%%%%%%%%%%%%%%%%%%%

The number $\astar$ now appears in three proved statements: the KPZ curve of the APPS model (as
$1+x_2$), the near-split lower bound, and the glued disks. None concerns finite meanders. This section compares
the missing step with percolation, where the corresponding step was carried out, and proves what the
missing pieces would give.

\subsection{The analogy with percolation}

In site percolation on the triangular lattice, each site is open independently with probability
$1/2$. The probability that the open cluster of a site, the sites it reaches through open
neighbors, reaches distance $R$ is $R^{-5/48+o(1)}$, proved in three pieces. \emph{The continuum
exponent}: Lawler, Schramm and Werner computed it for the continuum limit, built from SLE$_6$, the
curve describing boundaries of large clusters~\cite{LSW}. \emph{Convergence}: Smirnov proved
convergence to that limit~\cite{Smirnov}. \emph{Transfer of rare events}: convergence says nothing
about small probabilities, and Kesten's scaling relations, based on splitting a rare event into
nearly independent events on nested annuli, transfer the continuum probability to the
lattice~\cite{Kesten}. Together these give the lattice theorem~\cite{LSW,SW}; no single piece
suffices.

For meanders the continuum exponent is partly done: Theorem~\ref{thm:nearsplit} has the predicted
exponent on the lower side and $\xstar+5/7\approx3.1344$ on the upper side. Convergence is open:
it is Conjecture~\ref{conj:bgs}, and even the weaker Conjecture~\ref{conj:permuton} is open. The
transfer of rare events has barely started: a meander is encoded by its two systems of arches,
above and below the river, and this encoding has no independence between scales, so there is no
analog of nested annuli. What is known is combinatorial.

\subsection{Combinatorial transfer facts}

Write $L(n)=n\log\mu-\log M_n$, which is nonnegative because $\mu=\sup_nM_n^{1/n}$, and
$D(m)=\log(M_{2m}/M_m^2)=-\log q_m$. A power law $M_n\approx A\mu^nn^{-\alpha}$ means
$L(n)\approx\alpha\log n$.

\begin{lemma}\label{lem:transfer}
The following hold for all $m,n\ge1$.
\begin{enumerate}
\item[(a)] \emph{The halving identity.} $L(n)=\sum_{j\ge0}2^{-(j+1)}D(2^jn)$.
\item[(b)] \emph{Exact neck at one-cell tolerance.} For the centered atomic plot of a meander of
order $2m$ and $0<\delta<1/(8m)$, the near-split event $E_\delta$ is the zero-defect event $F_0$,
and $2q_m\le\Prob_{2m}(E_\delta)\le\frac12C_r^2q_m$.
\item[(c)] \emph{Few mismatches cost little.} There is an absolute constant $C$ such that the number
of meanders of order $m+n$ with at most $w\ge2$ road arches crossing a given river gap, with $2m$
crossings on its left, is at most $e^{Cw}M_mM_n$.
\item[(d)] \emph{Near-splits at small tolerance.} Consequently, for the centered atomic plot of a
uniform meander of order $2m$ and every $0<\delta<1/8$,
\[
q_m\le\Prob_{2m}(E_\delta)\le e^{C(m\delta+1)}q_m .
\]
\end{enumerate}
\end{lemma}

Part (a) holds because $D(m)=2L(m)-L(2m)$, so the sum telescopes, and $L(n)/n\to0$. Part (b) is
Proposition~\ref{prop:micro} with~\eqref{eq:F0}. Part (c) is a completion argument: close each half
into a meander by a few new crossings, merging the resulting loops one by one, and record the moves
in a word of length $w-1$ over four letters from which the construction can be undone. Part (d) follows from (c), since a near-split
with tolerance $\delta$ has of order $m\delta+1$ arches over the middle gap, and because an exact
junction is a near-split. So at tolerance $w/m$ with $w=o(\log m)$, a near-split costs the exact
neck probability times $m^{o(1)}$; the continuum sees much coarser tolerances.

\subsection{The conditional lower bound}

\begin{proposition}[conditional lower bound]\label{prop:condlower}
Suppose that for some $w_n\ge1$ with $w_n/\log n\to0$, the centered atomic plot $\Pi_n$ of a
uniform meander of order $n$ satisfies
\[
\Prob\bigl(\Pi_n\in E_{w_n/n}\bigr)\;\ge\;n^{-o(1)}\;\Prob\bigl(\Pi\in E_{w_n/n}\bigr)
\]
along even $n$. Then $M_n\ge\mu^nn^{-\astar-o(1)}$ for all $n$.
\end{proposition}

\begin{proof}
Let $n=2m$ and $\delta=w_n/n$. By Lemma~\ref{lem:transfer}(d), the hypothesis and the lower bound of
Theorem~\ref{thm:nearsplit},
\[
q_m\ge e^{-C(w_n+1)}\Prob(\Pi_n\in E_\delta)\ge e^{-C(w_n+1)}n^{-o(1)}c_1\Bigl(\frac{w_n}n\Bigr)^{\astar}
=n^{-\astar-o(1)},
\]
because $w_n=o(\log n)$ and $w_n\ge1$. So $D(m)\le(\astar+o(1))\log m$ as $m\to\infty$. Insert
this into the halving identity. Every argument $2^jn$ there is at least $n$, so for large $n$ the
bound applies to every term, and since $\sum_j2^{-(j+1)}\log(2^jn)=\log n+\log2$,
\[
L(n)\le(\astar+o(1))(\log n+\log2).
\]
Hence $M_n=\mu^ne^{-L(n)}\ge\mu^nn^{-\astar-o(1)}$.
\end{proof}

The hypothesis is a \emph{relative} estimate: it compares the probabilities of polynomially rare
events, at a tolerance just above one lattice cell, for finite meanders and for the limit object. It
is much stronger than Conjecture~\ref{conj:permuton}, which concerns only typical behavior. The conclusion would be the first power-law lower
bound for meanders; the best unconditional one, $M_n\ge\mu^ne^{-K\sqrt n}$, is in Paper~C~\cite{PaperC}.
% Checked: Paper C, Theorem 1, states this bound.

\subsection{The conditional upper bound}

The other direction needs much less precision. A \emph{coupling} of two random objects is a
construction of both on one probability space. For two measures $\rho,\rho'$ on $[0,1]^2$ let
$d_{\mathrm{rect}}(\rho,\rho')=\sup|\rho(\mathcal A\times\mathcal B)-\rho'(\mathcal A\times\mathcal B)|$, the supremum over rectangles
$\mathcal A\times\mathcal B$ with $\mathcal A$ an interval of river positions and $\mathcal B$ an interval of
road positions.

\begin{proposition}[conditional upper bound]\label{prop:condupper}
Suppose that for some $\sigma,\tau>0$ and every large even $n$, the centered atomic plot $\Pi_n$ of a
uniform meander of order $n$ and the meandric permuton $\Pi$ can be built on one probability space
so that
\[
\Prob\bigl(d_{\mathrm{rect}}(\Pi_n,\Pi)>n^{-\tau}\bigr)\le Cn^{-\sigma}.
\]
Then $M_n\le C_{\alpha'}\,\mu^nn^{-\alpha'}$ for every $\alpha'<\min\bigl(\sigma,(\xstar+5/7)\,\tau\bigr)$, where
$(\xstar+5/7)\tau\approx3.1344\,\tau$.
\end{proposition}

\begin{proof}
The off-block mass $d_\rho(s)$ is a sum of the masses of at most four rectangles, so
$|d_\rho(s)-d_{\rho'}(s)|\le4d_{\mathrm{rect}}(\rho,\rho')$ for every $s$. An exact junction of an order
$n=2m$ meander has $d_{\Pi_n}(s)=0$ for a suitable $s$. So on the event of a good coupling,
$\inf_sd_\Pi(s)\le4n^{-\tau}$, and by~\eqref{eq:sandwich} and Theorem~\ref{thm:nearsplit},
\[
q_m\le Cn^{-\sigma}+\Prob\bigl(\Pi\in E_{4n^{-\tau}}\bigr)\le Cn^{-\sigma}+C'n^{-(\xstar+5/7)\tau}.
\]
Hence $D(m)\ge\alpha'\log m-O(1)$ for every $\alpha'$ below the minimum, and the halving identity
gives $L(n)\ge\alpha'\log n-O(1)$ for all large $n$.
\end{proof}

Here the tolerance $n^{-\tau}$ may be much coarser than one cell, and no completion is needed. Since
$M_n\le\mu^{n+1}/(8n)$ is known~\cite{PaperC}, this improves on known bounds only when
$\min(\sigma,3.1344\,\tau)$ exceeds one.
% Checked: Paper C, Theorem 1, states this bound.
